
This section analyzes the economic properties of the protocol, including actor incentives, equilibrium conditions, and behavior under stress.

\paragraph{Scope and assumptions.} All results in this section are conditional on the external assumptions stated in Chapter 3: oracle liveness, underlying asset solvency, and bounded transaction costs. Claims about equilibrium behavior assume competitive arbitrage and sufficient market liquidity.

\section{Actors and Incentives}

\subsection{PT Holders (Fixed-Rate Exposure)}

\paragraph{Decision rule.}
A rational actor buys PT when the market-implied fixed rate exceeds their required return:
\[
r_{\mathrm{implied}} \geq r_{\mathrm{alt}} + \pi
\]
where $r_{\mathrm{alt}}$ is the opportunity cost of capital and $\pi$ is a risk premium covering contract risk, oracle risk, and underlying impairment risk.

\paragraph{Implied rate.}
Let $\Delta$ be time-to-expiry in years and $P_{PT}$ the PT price in SY units. The implied annualized return is:
\[
r_{\mathrm{implied}} = P_{PT}^{-1/\Delta} - 1
\]
For example, if $P_{PT} = 0.95$ and $\Delta = 0.5$ years, the implied APY is $(0.95)^{-2} - 1 \approx 10.8\%$.

\paragraph{Risks.}
\begin{itemize}
\item \textbf{Smart contract risk}: Low probability, high severity. Mitigated by audits and formal verification.
\item \textbf{Underlying impairment}: PT redemption is deterministic in SY units, but SY value depends on underlying solvency.
\item \textbf{Opportunity cost}: Capital is locked; holder cannot chase higher yields elsewhere.
\item \textbf{Early exit}: Must sell on AMM at prevailing price, subject to slippage and mark-to-market loss.
\end{itemize}

\paragraph{Implications.}
PT provides deterministic redemption in SY units at expiry, conditional on SY solvency. Prior to expiry, PT is mark-to-market and may trade at a discount or premium depending on rate expectations and liquidity.

\subsection{YT Holders (Yield Speculation)}

\paragraph{Decision rule.}
A rational actor buys YT when expected yield exceeds the YT price:
\[
\mathbb{E}[g(\Delta)] \geq P_{YT}
\]
where $g(\Delta)$ is the expected fractional yield growth of the underlying over time-to-expiry $\Delta$. Equivalently, if $y$ is the expected APY:
\[
1 - e^{-y\Delta} \geq P_{YT}
\]
This directly ties to the conservation identity $P_{PT} + P_{YT} = 1$.

\paragraph{Leverage effect.}
YT provides leveraged exposure to yield. If $P_{YT} = 0.05$ (5\% of SY value) and the underlying yields 10\% APY over a 6-month period:
\begin{itemize}
\item SY holder earns: $\approx 5\%$ on 100\% capital
\item YT holder earns: $\approx 5\%$ yield on 5\% capital = $\approx 100\%$ effective return on YT investment
\end{itemize}

\paragraph{Risks.}
\begin{itemize}
\item \textbf{Time decay}: YT value converges to zero as expiry approaches, regardless of yield.
\item \textbf{Yield drops}: Reduced interest collection; YT loses value rapidly.
\item \textbf{Complete depeg}: If underlying stops generating yield, YT becomes worthless.
\end{itemize}

\paragraph{Implications.}
YT is a high-risk, high-reward instrument suitable for actors with strong convictions about future yield. Longer maturities provide more leverage but greater time decay exposure.

\subsection{Liquidity Providers}

\paragraph{Decision rule.}
LP provision is attractive when expected return exceeds opportunity cost:
\[
\mathbb{E}[\text{Fee APY}] + r_{SY} \cdot w_{SY} - \mathbb{E}[\text{IL}] \geq r_{\mathrm{alt}}
\]
where $r_{SY}$ is the SY yield rate, $w_{SY}$ is the SY weight in the pool, and IL is impermanent loss.

\paragraph{Impermanent loss dynamics.}
Unlike traditional AMMs where IL is symmetric, PT/SY pools exhibit asymmetric and time-dependent IL:

\begin{center}
\begin{tabular}{lll}
\toprule
\textbf{Rate Movement} & \textbf{PT Price} & \textbf{LP Impact} \\
\midrule
Rates increase & PT decreases & IL: LP holds excess PT \\
Rates decrease & PT increases & IL: LP holds excess SY \\
Near expiry & PT $\to$ 1 & IL minimizes as prices converge \\
\bottomrule
\end{tabular}
\end{center}

\paragraph{Time-to-expiry effect.}
As expiry approaches, the PT price becomes less sensitive to rate changes (via the rate scalar $r_{scalar} \to \infty$), reducing volatility-driven IL. However, this also typically reduces trading volume and fee opportunities. LPs face a tradeoff between IL protection and fee income as maturity nears.

\paragraph{Risks.}
\begin{itemize}
\item \textbf{Directional IL}: Unlike xy=k pools, rate movements cause directional exposure shifts.
\item \textbf{Rate volatility}: High volatility increases IL without proportional fee compensation.
\item \textbf{Low volume periods}: Fee income may not offset capital opportunity cost.
\end{itemize}

\paragraph{Implications.}
LP returns are path-dependent and sensitive to rate volatility. Optimal entry is during stable rate environments; exit timing around volatility events is critical.

\subsection{Arbitrageurs}

\paragraph{Conservation bounds.}
The conservation identity creates soft bounds on prices:
\[
P_{PT} + P_{YT} = 1 \pm \epsilon
\]
where $\epsilon$ is bounded by transaction costs, slippage, and execution latency. In liquid markets with low gas costs, $\epsilon < 0.5\%$ is typical.

\paragraph{Arbitrage mechanisms.}

\begin{center}
\begin{tabular}{lll}
\toprule
\textbf{Condition} & \textbf{Action} & \textbf{Effect} \\
\midrule
$P_{PT} + P_{YT} > 1 + \epsilon$ & Mint PT+YT, sell both & Prices decrease \\
$P_{PT} + P_{YT} < 1 - \epsilon$ & Buy PT+YT, redeem for SY & Prices increase \\
Horizon yield $>$ external & Sell PT (or buy YT) & Implied yield decreases \\
Horizon yield $<$ external & Buy PT (or sell YT) & Implied yield increases \\
\bottomrule
\end{tabular}
\end{center}

\paragraph{Implications.}
Arbitrage activity maintains price consistency with external rate expectations and enforces the conservation identity within transaction cost bounds. Arbitrage profitability decreases as markets become more efficient.

\section{Equilibrium Conditions}

\subsection{Price Equilibrium}

At equilibrium, the implied yield reflects market consensus on future rates:
\[
r_{\mathrm{implied}} = \mathbb{E}[r_{\mathrm{actual}}] + \pi
\]
where $\pi$ is a risk premium for contract and oracle risk. Deviations from this condition create arbitrage opportunities that are corrected by rational traders.

\subsection{Liquidity Equilibrium}

LP provision reaches equilibrium when marginal return equals opportunity cost:
\[
\text{Fee APY} + r_{SY} \cdot w_{SY} - \mathbb{E}[\text{IL}] = r_{\mathrm{alt}}
\]

Factors affecting LP equilibrium depth:
\begin{itemize}
\item Higher trading volume $\Rightarrow$ more fees $\Rightarrow$ increased LP supply
\item Higher rate volatility $\Rightarrow$ more IL $\Rightarrow$ decreased LP supply
\item Higher underlying yield $\Rightarrow$ more attractive SY exposure $\Rightarrow$ increased LP supply
\end{itemize}

\section{Stress Scenarios}

\subsection{Yield Collapse}

\paragraph{Scenario.} Underlying yield drops from 10\% to 2\% APY.

\begin{center}
\begin{tabular}{lll}
\toprule
\textbf{Actor} & \textbf{Immediate Impact} & \textbf{Final Outcome} \\
\midrule
PT Holder & Mark-to-market decline & Converges to 1 SY at maturity \\
YT Holder & YT price drops 80\%+ & Collects only 2\% yield \\
LP & IL + reduced fee volume & Depends on exit timing \\
\bottomrule
\end{tabular}
\end{center}

\paragraph{Key observation.} PT holders who hold to maturity are unaffected in SY terms; the ``loss'' is opportunity cost relative to new higher-yielding alternatives.

\subsection{Yield Spike}

\paragraph{Scenario.} Underlying yield spikes from 10\% to 25\% APY.

\begin{center}
\begin{tabular}{lll}
\toprule
\textbf{Actor} & \textbf{Immediate Impact} & \textbf{Final Outcome} \\
\midrule
PT Holder & Opportunity cost (missed yield) & Still receives 1 SY at maturity \\
YT Holder & Large windfall & Collects 25\% yield; YT appreciates \\
LP & Directional IL & Higher volume partially offsets \\
\bottomrule
\end{tabular}
\end{center}

\subsection{Mass Redemption}

\paragraph{Scenario.} 80\% of PT+YT holders redeem simultaneously.

\paragraph{Key property.} Unlike lending protocols, the protocol has no maturity mismatch or run-induced insolvency risk, assuming SY is fully collateralized by the underlying. Each PT+YT pair is backed by exactly 1 SY held in the YT contract; redemption is always possible.

\paragraph{Constraints.}
\begin{itemize}
\item Redemption throughput is limited by on-chain block capacity.
\item AMM exits are subject to price impact and slippage.
\item Solvency does not depend on borrower repayment (unlike lending protocols).
\end{itemize}

\subsection{Underlying Depeg}

\paragraph{Scenario.} Underlying token loses 30\% of external value.

\begin{center}
\begin{tabular}{ll}
\toprule
\textbf{Asset} & \textbf{Impact} \\
\midrule
SY & Worth 0.7$\times$ previous external value \\
PT & Redeems for 1 SY (now worth 0.7$\times$) \\
YT & Yield continues at depressed rate \\
LP & Proportional loss on both reserve assets \\
\bottomrule
\end{tabular}
\end{center}

\paragraph{Key distinction.} Depeg affects the value of SY in external terms (e.g., USD), not the internal accounting identity PT $\to$ 1 SY. The protocol's internal guarantees remain intact; external value loss is passed through to all token holders.

\paragraph{Protocol response.}
\begin{itemize}
\item Watermark index prevents distributing negative yield.
\item All operations continue functioning normally.
\item Users can exit at prevailing market prices.
\end{itemize}

\section{Mechanism Properties}

\subsection{Equilibrium Characterization}

The system admits a candidate equilibrium where no actor can improve their expected payoff by unilaterally deviating:

\begin{center}
\begin{tabular}{ll}
\toprule
\textbf{Actor} & \textbf{Equilibrium Condition} \\
\midrule
PT Buyer & Implied yield = required return for risk \\
YT Buyer & Expected yield $\geq$ YT price + speculation premium \\
LP & Fee income = IL + opportunity cost \\
Arbitrageur & Spread $\leq$ execution costs \\
\bottomrule
\end{tabular}
\end{center}

Under competitive arbitrage and sufficient liquidity, deviations from equilibrium prices are mean-reverting.

\subsection{Incentive Properties}

\paragraph{Incentive compatibility.}
Honest behavior (accurate price discovery through trading) is locally incentive-compatible under competitive arbitrage and bounded manipulation costs. Price manipulation is costly (requires capital) and temporary (arbitrage corrects).

\paragraph{Individual rationality.}
Participation is voluntary. Expected utility is non-negative for informed participants who correctly assess risks and opportunities.

\paragraph{Budget balance.}
The protocol is self-sustaining: fees cover operational costs (gas for oracle updates, etc.). No external subsidies are required for equilibrium operation.

\section{Fee Economics}

\subsection{Fee Retention Model}

Fees are retained in pool reserves rather than extracted:
\begin{itemize}
\item Swap fees increase reserve balances (SY or PT depending on trade direction).
\item LP token value appreciates proportionally to accumulated fees.
\item Fee value is realized when LPs withdraw liquidity.
\end{itemize}

\paragraph{Terminology note.} The swap mechanism uses a \texttt{fee\_multiplier} $\in [0,1]$ where:
\[
\text{fee} = \text{gross\_amount} \times (1 - \text{fee\_multiplier})
\]
A fee multiplier of 0.999 corresponds to a 0.1\% fee rate.

\subsection{Optimal Fee Rate}

The optimal fee rate balances trader demand and LP incentives:
\begin{itemize}
\item \textbf{Too low}: Insufficient LP incentive $\Rightarrow$ thin liquidity $\Rightarrow$ high slippage.
\item \textbf{Too high}: Traders use alternatives $\Rightarrow$ low volume $\Rightarrow$ reduced LP returns.
\end{itemize}

Market forces determine the effective optimum:
\begin{itemize}
\item Competition between markets for the same underlying sets fee ceiling.
\item Arbitrage activity requires fees below arbitrage profit margins.
\item LP entry/exit adjusts liquidity depth to equilibrium levels.
\end{itemize}

\section{Value Accrual}

\subsection{Value Flow}

\begin{center}
\begin{tabular}{ll}
\toprule
\textbf{Value Source} & \textbf{Recipient} \\
\midrule
Underlying yield & YT holders (until expiry) \\
PT discount capture & PT holders (at maturity) \\
Swap fees & LP holders \\
Arbitrage spreads & Arbitrageurs \\
\bottomrule
\end{tabular}
\end{center}

\paragraph{Zero-sum framing.}
Discount capture (PT) and yield capture (YT) are two sides of the same decomposition: PT holders forgo yield in exchange for certainty; YT holders accept uncertainty for leveraged yield exposure. Fees and arbitrage spreads are paid by traders through execution price.

\subsection{Protocol Value Capture}

\paragraph{Current design.} Zero protocol fee. All value flows to users (PT/YT/LP holders and arbitrageurs).

\paragraph{Future consideration.} A governance-controlled protocol fee on swaps could direct value to a treasury or token holders. This would be implemented as a fraction of swap fees retained by the protocol rather than accruing to LPs.

